\section{Results}
\label{sec:results}

This section presents the main results of our analysis, including comparative statics, welfare analysis, and policy implications. All results use the calibrated parameters from Section \ref{sec:calibration}.

\subsection{Comparative Statics}

\subsubsection{Herding Strength}

Figure \ref{fig:herding} shows the effects of varying herding strength ($\alpha_{\text{herd}}$) from 0.01 to 0.20.

\begin{figure}[ht]
\centering
\includegraphics[width=\textwidth]{figures/figure2_herding_comparative.pdf}
\caption{Comparative Statics: Herding Strength Effects}
\label{fig:herding}
\end{figure}

\textbf{Key findings:}

\begin{enumerate}
    \item \textbf{Kurtosis}: Increases from 3.2 to 7.8 (+144\%) as herding increases. This confirms that herding behavior generates fat tails, consistent with \citet{shiller1995conversation}.
    
    \item \textbf{Welfare}: Decreases from -0.82 to -0.91 (-11\%), indicating efficiency losses from excessive herding.
    
    \item \textbf{Crash risk}: Increases from 0.02 to 0.06 (+200\%), suggesting that herding amplifies systemic risk.
    
    \item \textbf{Mispricing}: Increases from 0.08 to 0.15 (+88\%), as herding drives prices away from fundamentals.
\end{enumerate}

\textbf{Policy implication:} Moderate herding may be beneficial (maintains diversity), but excessive herding should be monitored and potentially constrained.

\subsubsection{Learning Speed}

Figure \ref{fig:learning} shows the effects of varying learning speed ($\phi$) from 0.60 to 0.99.

\begin{figure}[ht]
\centering
\includegraphics[width=\textwidth]{figures/figure5_robustness.pdf}
\caption{Comparative Statics: Learning Speed Effects}
\label{fig:learning}
\end{figure}

\textbf{Key findings:}

\begin{enumerate}
    \item \textbf{Volatility}: Decreases from 0.023 to 0.018 (-22\%) as learning speeds up. Faster learning leads to quicker convergence and lower volatility.
    
    \item \textbf{Efficiency}: Increases from 0.89 to 0.94 (+6\%), consistent with the intuition that better learning improves price discovery.
    
    \item \textbf{Diversity}: Decreases from 0.82 to 0.71 (-13\%), as fast learning may lead to herding on popular strategies.
\end{enumerate}

\textbf{Policy implication:} Promoting information dissemination (e.g., through transparency requirements) can improve market efficiency, but policymakers should be aware of potential diversity losses.

\subsubsection{Selection Sensitivity}

Figure \ref{fig:selection} shows the effects of varying selection sensitivity ($\lambda$) from 0.5 to 5.0.

\begin{figure}[ht]
\centering
\includegraphics[width=\textwidth]{figures/figure4_tradeoff.pdf}
\caption{Comparative Statics: Selection Sensitivity Effects}
\label{fig:selection}
\end{figure}

\textbf{Key findings:}

\begin{enumerate}
    \item \textbf{Inequality}: Increases from 0.72 to 0.85 (+18\%), as high selection leads to ``winner-takes-all'' dynamics.
    
    \item \textbf{Mispricing}: Decreases from 0.08 to 0.03 (-63\%), as strong selection eliminates inefficient strategies.
\end{enumerate}

\textbf{Policy implication:} There is a trade-off between efficiency and equality. Policymakers need to balance these competing objectives.

\subsubsection{Leverage Limits}

Figure \ref{fig:leverage} shows the effects of varying leverage limits from 2x to 50x.

\begin{figure}[ht]
\centering
\includegraphics[width=\textwidth]{figures/figure3_welfare_policy.pdf}
\caption{Comparative Statics: Leverage Limit Effects}
\label{fig:leverage}
\end{figure}

\textbf{Key findings:}

\begin{enumerate}
    \item \textbf{Efficiency}: Increases from 0.87 to 0.95 (+9\%) as leverage increases, since arbitrageurs can correct mispricing more effectively.
    
    \item \textbf{Mispricing}: Decreases from 0.12 to 0.05 (-58\%), confirming that leverage constraints limit arbitrage effectiveness.
    
    \item \textbf{Welfare}: Shows an inverted U-shape, peaking at 10x leverage. Too little leverage limits arbitrage; too much increases crisis risk.
\end{enumerate}

\textbf{Policy implication:} Moderate leverage limits (5--10x) balance efficiency and stability, consistent with \citet{shleifer1997limits}.

\subsection{Welfare Analysis}

\subsubsection{Social Welfare Function}

We define social welfare as a weighted utilitarian function:
\begin{equation}
\mathcal{W} = \sum_{k=1}^K \omega_k \cdot U_k
\end{equation}
where $U_k = -\exp(-\gamma W_k) - \frac{1}{2}\text{Var}(W_k)$ is the CARA utility of agent type $k$, and $\omega_k = 0.2$ are equal welfare weights.

Inequality is measured by the Gini coefficient:
\begin{equation}
G = \frac{\sum_{i=1}^n \sum_{j=1}^n |W_i - W_j|}{2 n^2 \cdot \bar{W}}
\end{equation}

\subsubsection{Policy Experiments}

We analyze seven regulatory policies:

\begin{enumerate}
    \item \textbf{Baseline}: Calibrated parameters
    \item \textbf{Short-sale ban}: Prohibit short selling ($\chi = 1$)
    \item \textbf{Leverage cap 3x}: Maximum leverage ratio of 3
    \item \textbf{Leverage cap 10x}: Maximum leverage ratio of 10
    \item \textbf{Transaction tax}: 0.1\% tax on all trades
    \item \textbf{High herding}: Increase herding strength to 0.15
    \item \textbf{Low learning}: Decrease learning speed to 0.6
\end{enumerate}

\subsubsection{Results}

Table \ref{tab:welfare} reports welfare effects for each policy.

\begin{table}[ht]
\centering
\caption{Welfare Effects of Regulatory Policies}
\label{tab:welfare}
\begin{tabular}{lrrr}
\hline\hline
Policy & Welfare Change (\%) & Gini Change & Ranking \\
\hline
Leverage cap 10x & +2.3 & -0.012 & 1 \\
Leverage cap 3x & +1.5 & -0.008 & 2 \\
Low learning & +0.8 & +0.003 & 3 \\
Baseline & 0.0 & -- & 4 \\
High herding & -0.5 & +0.015 & 5 \\
Transaction tax & -1.2 & +0.008 & 6 \\
Short-sale ban & -1.8 & +0.023 & 7 \\
\hline\hline
\end{tabular}
\end{table}

\textbf{Key findings:}

\begin{enumerate}
    \item \textbf{Best policy}: Leverage cap 10x increases welfare by 2.3\% and reduces inequality (Gini -0.012). This balances arbitrage effectiveness with financial stability.
    
    \item \textbf{Worst policy}: Short-sale ban reduces welfare by 1.8\% and increases inequality (Gini +0.023). This limits arbitrage and allows mispricing to persist.
    
    \item \textbf{Transaction tax}: Reduces welfare by 1.2\%, as it reduces liquidity and impairs price discovery.
    
    \item \textbf{High herding}: Reduces welfare by 0.5\%, confirming that excessive herding is harmful.
\end{enumerate}

\subsubsection{Welfare-Equality Trade-off}

Figure \ref{fig:tradeoff} shows the welfare-inequality trade-off across policies.

\begin{figure}[ht]
\centering
\includegraphics[width=\textwidth]{figures/figure4_tradeoff.pdf}
\caption{Welfare-Equality Trade-off}
\label{fig:tradeoff}
\end{figure}

Most policies show a trade-off: higher welfare comes with higher inequality. However, leverage caps (10x) are a ``win-win'' policy that improves both welfare and equality.

\subsection{Equilibrium Properties}

\subsubsection{Strategy Distribution}

Table \ref{tab:equilibrium} reports the equilibrium strategy distribution.

\begin{table}[ht]
\centering
\caption{Equilibrium Strategy Distribution}
\label{tab:equilibrium}
\begin{tabular}{lr}
\hline\hline
Strategy & Weight (\%) \\
\hline
Noise Trader & 20.45 \\
Momentum Trader & 18.67 \\
Value Trader & 31.22 \\
Herd Trader & 18.89 \\
Arbitrageur & 10.77 \\
\hline
\textbf{Total} & \textbf{100.00} \\
\hline\hline
\end{tabular}
\end{table}

\textbf{Key observations:}

\begin{enumerate}
    \item Value traders dominate (31.22\%), suggesting that the market is relatively efficient.
    
    \item Noise and herd traders together account for 39.34\%, indicating significant behavioral biases.
    
    \item Arbitrageurs are only 10.77\%, consistent with limits to arbitrage (\citealp{shleifer1997limits}).
    
    \item Diversity persists (no single strategy dominates), supporting H1 (heterogeneity).
\end{enumerate}

\subsubsection{Stability Verification}

We verify local stability by computing the Jacobian spectral radius:
\begin{itemize}
    \item Spectral radius: $\rho(J) = 0.7342 < 1$ \quad \textbf{✓ Stable}
    \item Maximum eigenvalue: 0.7342
    \item All eigenvalues inside unit circle
\end{itemize}

This confirms Proposition 2 (local stability) numerically.

\subsection{Discussion}

\subsubsection{Why These Results Matter}

Our results have three important implications:

\textbf{First}, from a \textit{theoretical} perspective, we show that agent-based models can be rigorous. Equilibrium existence and stability are not assumed---they are proved. This addresses a common criticism of ABM.

\textbf{Second}, from an \textit{empirical} perspective, we show that the model reproduces key stylized facts with economically meaningful parameters. This is not curve-fitting---the parameters are grounded in experimental literature.

\textbf{Third}, from a \textit{policy} perspective, we provide quantitative welfare analysis. Policymakers can use our framework to evaluate ``what-if'' scenarios that are impossible to test in the real world.

\subsubsection{Limitations}

Our results have limitations:

\begin{enumerate}
    \item \textbf{Single market}: We calibrate to S\&P 500 only. Cross-market validation would strengthen the results.
    
    \item \textbf{Simplified market microstructure}: We use a market maker model, not a full order book. This may miss some dynamics.
    
    \item \textbf{Parameter uncertainty}: While we use priors from the literature, there is uncertainty. Bayesian calibration would provide posterior distributions.
\end{enumerate}

Despite these limitations, the results are robust and economically meaningful.

\subsection{Conclusion}

This section presents the main results of our analysis. Comparative statics show how parameters affect market outcomes. Welfare analysis ranks seven regulatory policies, with leverage caps (10x) emerging as the optimal policy. Equilibrium analysis confirms that diversity persists and the system is stable.

These results provide strong support for the agent-based approach and demonstrate its practical relevance for policy analysis.
